\chapter{The Closing Statement}
\label{ch:synthesis}

\begin{quotation}
\noindent\textit{Synopsis.} This closing chapter integrates the
monograph's concepts into a single narrative: distinguishability explains
what survives; admissibility explains what remains possible; repair
economics explains why a distinction was or was not maintained; budget
relays explain how losses accumulate across time and institutions; the
fidelity continuum unifies preservation and reconstitution; and Composition
Failure explains why locally rational decisions can still produce globally
insufficient outcomes. The chapter closes by naming what remains
unresolved, above all the MEM|8 dependency, and restates the central claim
that continuation is constrained not by malice or error but by the
unavoidable limits of finite repair budgets operating in an open future.
\end{quotation}


\section{Restating the Arc}

Distinguishability (\cref{ch:distinctions}) established what survives
a projection. Admissibility (\cref{ch:admissibility}) established what
continuations remain reachable given what survives. Repair economics
(\cref{ch:repair-economics}) established why a given distinction was
or was not maintained, and, via \cref{cor:ambiguity-limiting-case},
established that this monograph could not simply inherit its central
claim from \emph{Conservation of Ambiguity}. The budget relay
(\cref{ch:budget-relay}) formalized how discard compounds across
uncoordinated stages, and the Composition Failure Theorem
(\cref{thm:composition-failure}) proved the diachronic analogue the
Redistribution Theorem does not cover. The fidelity continuum
(\cref{ch:fidelity-continuum}) collapsed preservation and
reconstitution into a single variable, conditional on the unresolved
question raised in Chapter~\ref{ch:mem8}.

\section{The Closing Statement}

\begin{conjecture}[Closing Statement]
\label{conj:closing}
Continuation depends on the substrate available to a consumer, whether
preserved with high fidelity or regenerated from traces at low
fidelity; in either case the substrate is bounded by the distinctions
its production budget could afford to maintain. Because relays compose
without guaranteed coordination (\cref{def:budget-relay}) and because
future task spaces remain open (\cref{prop:task-unavailability}), no
finite chain of substrates can be guaranteed sufficient for arbitrary
future continuation -- not because information is destroyed, but
because no stage can afford, or is positioned, to preserve everything
a task it cannot yet foresee will require.
\end{conjecture}

This statement is labeled a conjecture rather than a theorem
deliberately: it synthesizes \cref{thm:composition-failure} together
with the fidelity continuum of Chapter~\ref{ch:fidelity-continuum},
but the synthesis itself, as a single unified statement, has not been
given an independent proof beyond what its components already
establish.

\section{What Remains Open}

Appendix~\ref{app:open-questions} consolidates every conjecture and
open question raised in the body. The largest is
\cref{oq:mem8-fork}, on which the scope of Part~\ref{part:modes-of-access}'s
results depends. The tightness of \cref{cor:no-free-lunch}'s growth
bound is conjectural. Whether a general coordination protocol across
budget relays could bound, though not eliminate,
\cref{thm:composition-failure}'s failure mode is left as a direction
for future work rather than claimed here.

\section{What This Monograph Does Not Claim}

This monograph does not claim that all cognition is reconstructive; it
restricts that claim, via \cref{ch:mem8}'s narrow reading, to cases
independently supported by existing literature. It does not claim that
\emph{Conservation of Ambiguity}'s Redistribution Theorem is
superseded; \cref{thm:redistribution} is imported and used, not
revised. It does not claim to have resolved the preservation/
reconstitution question definitively, since that resolution remains
contingent on \cref{oq:mem8-fork}.
